%%=============================================================================
%% Methodologie
%%=============================================================================

\chapter{\IfLanguageName{dutch}{Methodologie}{Methodology}}%
\label{ch:methodologie}

%% TODO: In dit hoofstuk geef je een korte toelichting over hoe je te werk bent
%% gegaan. Verdeel je onderzoek in grote fasen, en licht in elke fase toe wat
%% de doelstelling was, welke deliverables daar uit gekomen zijn, en welke
%% onderzoeksmethoden je daarbij toegepast hebt. Verantwoord waarom je
%% op deze manier te werk gegaan bent.
%% 
%% Voorbeelden van zulke fasen zijn: literatuurstudie, opstellen van een
%% requirements-analyse, opstellen long-list (bij vergelijkende studie),
%% selectie van geschikte tools (bij vergelijkende studie, "short-list"),
%% opzetten testopstelling/PoC, uitvoeren testen en verzamelen
%% van resultaten, analyse van resultaten, ...
%%
%% !!!!! LET OP !!!!!
%%
%% Het is uitdrukkelijk NIET de bedoeling dat je het grootste deel van de corpus
%% van je bachelorproef in dit hoofstuk verwerkt! Dit hoofdstuk is eerder een
%% kort overzicht van je plan van aanpak.
%%
%% Maak voor elke fase (behalve het literatuuronderzoek) een NIEUW HOOFDSTUK aan
%% en geef het een gepaste titel.

%%=============================================================================
%% Methodologie
%%=============================================================================

\section{Methodologie}

In deze bachelorproef wordt een toegepaste kwalitatieve case study uitgevoerd bij VLAIO naar de meerwaarde van OpenMetadata voor de transparantie, herleidbaarheid en ervaren begrijpbaarheid van rapporten en impactanalyses. De nadruk ligt daarbij niet op een strikte voor- en nameting met vaste procesindicatoren, maar op een praktijkgerichte evaluatie van een proof-of-concept (PoC) en de manier waarop deze door betrokken gebruikers ervaren en beoordeeld wordt.\\

Deze keuze sluit aan bij de specifieke context van VLAIO. Wanneer binnen de dataomgeving een wijziging gepland wordt, bijvoorbeeld aan een databron, datamodel of verwerking, worden soms technische impactanalyses opgesteld om in kaart te brengen welke datasets, tabellen en rapporten daardoor geraakt kunnen worden en welke acties ondernomen moeten worden om de wijziging succesvol tot stand te brengen. Dergelijke impactanalysedocumenten worden in de praktijk echter slechts door een beperkte groep technische profielen opgesteld, zoals data-engineers en analisten. Voor deze groep ligt de verwachte meerwaarde van lineage bovendien relatief voor de hand, aangezien een expliciet overzicht van afhankelijkheden het eenvoudiger maakt om sneller betrokken datasets en rapporten te identificeren en wijzigingen beter te onderbouwen.\\

Tegelijk is dit technische impactanalyseproces minder geschikt voor een gestandaardiseerde kwantitatieve evaluatie. De inhoud, complexiteit en aanleiding van impactanalyses variëren sterk, waardoor indicatoren zoals doorlooptijd of aantal geïdentificeerde afhankelijkheden moeilijk op een betrouwbare en vergelijkbare manier te meten zijn. Impactanalyses zijn immers geen strikt afgelijnde, repetitieve taken waarbij je een verbetering kan testen met een voor en nameting.\\

Daarnaast wordt impactanalyse in deze bachelorproef breder opgevat dan enkel het formeel opstellen van zulke technische impactanalyse documenten. Ook het begrijpen van rapporten, het duiden van wijzigingen en het inschatten van afhankelijkheden voor een ruimere groep interne gebruikers maken deel uit van die bredere context.\\

Om die redenen wordt de evaluatie in deze bachelorproef niet primair gericht op het technische impactanalyseproces zelf, maar op de bredere manier waarop inzicht in dataherkomst gebruikers ondersteunt bij het begrijpen van rapporten, het interpreteren van wijzigingen en het inschatten van afhankelijkheden. Deze benadering maakt het mogelijk om de meerwaarde van lineage te onderzoeken voor een ruimere en diversere doelgroep van rapportgebruikers, rapportbouwers en technische profielen, waar de ervaren meerwaarde minder vanzelfsprekend is en daarom relevanter is om empirisch te onderzoeken.\\

De methodologie bestaat uit drie fasen: 1) analyse en afbakening van de praktijkcontext, 2) onderzoek en implementatie van een proof-of-concept in OpenMetadata, en 3) evaluatie en interpretatie van de proof-of-concept op basis van gebruikersfeedback. Deze gefaseerde aanpak laat toe om het onderzoek systematisch op te bouwen, de technische realisatie van lineage te koppelen aan concrete praktijkcases, en de ervaren meerwaarde, beperkingen en aandachtspunten op een onderbouwde manier in kaart te brengen.\\

\subsection{Fase 1: Analyse van de praktijkcontext en afbakening van de scope}

Het doel van deze fase is om de huidige werkwijze rond impactanalyses, rapportering en zicht op dataherkomst binnen VLAIO beter te begrijpen en om de scope van het praktijkonderzoek af te bakenen. Hiervoor worden bestaande impactanalyse-documenten, relevante rapporteringsketens, de data architectuur en de betrokken dataflows verkend. In Hoofdstuk~\ref{ch:contextanalyse} wordt deze analyse samengebracht in een beschrijving van de Databricks-lakehouseomgeving, de medallion-architectuur, de huidige mogelijkheden en beperkingen van lineage, en de manier waarop impactanalyses vandaag verlopen.\\

Op basis van deze contextanalyse wordt de scope van de proof-of-concept bepaald. Daarbij ligt de focus op het generiek zichtbaar maken van de ontbrekende koppeling tussen platinum-datasets in Databricks en Power BI-rapporten in OpenMetadata. De onderliggende aanpak is dus niet ontwikkeld voor één specifiek rapport of domein, maar met het oog op bredere toepasbaarheid binnen het Power BI-landschap van VLAIO. Voor het testen en valideren van de werking wordt in de implementatiefase wel gebruikgemaakt van een beperkte subset van rapporten, louter als praktische testomgeving om de proof-of-concept stapsgewijs op te bouwen en te controleren.\\

Daarnaast wordt in deze fase bepaald welke profielen in het evaluerende deel betrokken worden. Daarbij ligt de focus op rapportgebruikers, rapportbouwers en technische dataprofielen, omdat de evaluatie niet alleen peilt naar de ondersteuning van impactanalyses, maar ook naar de begrijpbaarheid van rapporten en de interpretatie van wijzigingen. Op die manier wordt de evaluatie afgestemd op de bredere onderzoeksdoelstelling van deze bachelorproef. Voor de bevraging wordt een bredere groep van ongeveer 80 interne medewerkers aangesproken, met de verwachting dat dit voldoende respons oplevert om minstens een twintigtal bruikbare antwoorden te verzamelen.\\

\textbf{Deliverable}: een afgebakende praktijkcontext met beschrijving van de data-archi\-tec\-tuur, de huidige lineage-lacunes richting Power BI, een duidelijke scope voor een generiek inzetbare proof-of-concept, en een lijst van relevante personen voor de evaluatie.

\subsection{Fase 2: Ontwerp en implementatie van de proof-of-concept}

In deze fase wordt gezocht naar een haalbare oplossing voor de ontbrekende zichtbaarheid tussen platinum-datasets in Databricks en Power BI-rapporten. Daarbij wordt OpenMetadata geconfigureerd in functie van de beschikbare metadata over Databricks-tabellen en Power BI-rapporten, en wordt een generieke naamgebaseerde mappinglaag opgebouwd die Power BI-datasetnamen koppelt aan onderliggende platinum-datasets.\\

De implementatie maakt gebruik van Databricks data uit de Power BI REST API en een export van platinum-datasets, die in scripts en notebooks worden ingelezen en verwerkt. De mappinggenerator wordt toegepast op de volledige set Power BI-datasets uit de export en levert voor een aanzienlijk deel daarvan bruikbare mappings naar platinum-datasets op. De volledige end-to-end lineage-flow, inclusief dashboardcreatie en het effectief leggen van lineage-edges, wordt vervolgens stapsgewijs getest en gevalideerd in een beperkte praktische testopstelling. De onderliggende scripts en YAML-configuraties zijn echter generiek opgezet, zodat dezelfde aanpak nadien ook breder inzetbaar is voor andere rapporten en datasets.\\

De proof-of-concept wordt uitgewerkt als een combinatie van herbruikbare Pythonscripts, Jupyter-notebooks en YAML-configuraties. De nadruk ligt daarbij op het inlezen en verwerken van beschikbare Power BI- en Databricks-metadata, het genereren van mappings tussen Power BI-datasets en platinum-datasets, en het programmatisch aanmaken van dashboards en lineage-edges in OpenMetadata. Daarnaast worden de belangrijkste ontwerpkeuzes en beperkingen vastgelegd, zoals de keuze voor lineage op tabelniveau, de afhankelijkheid van naamgebaseerde mapping en de beperkte mogelijkheid om de end-to-end-flow onmiddellijk voor het volledige rapportagelandschap uit te voeren.\\

\textbf{Deliverable}: een werkende proof-of-concept met OpenMetadata, bestaande uit herbruikbare scripts, notebooks en YAML-configuraties voor het zichtbaar maken van de koppeling tussen Power BI-rapporten en onderliggende platinum-datasets, inclusief documentatie van ontwerpkeuzes, technische bevindingen en beperkingen.

\subsection{Fase 3: Evaluatie en interpretatie van de proof-of-concept}

In de derde fase wordt de proof-of-concept geëvalueerd op haar bruikbaarheid en ervaren meerwaarde voor impactanalyses en rapportgebruik binnen VLAIO. De kern van deze fase is de online vragenlijst \emph{``Begrijpbaarheid en vertrouwen in rapporten''}, opgesteld in Microsoft Forms en gericht op gebruikers die in meer of mindere mate met rapporten, zoals Power BI, werken. De vragenlijst bevat profielvragen, stellingen over de huidige ervaring zonder extra inzicht in dataherkomst, een scenario met een herkomstoverzicht gebaseerd op de gerealiseerde PoC, bijkomende vragen voor technische rapportbouwers en dataprofielen, en een algemene evaluatie van de meerwaarde en het verwachte gebruik van zo'n overzicht. De meeste stellingen worden beantwoord op een vijfpuntsschaal van ``helemaal oneens'' tot ``helemaal eens'', aangevuld met open vragen waarin respondenten concrete situaties, barrières of suggesties kunnen beschrijven.\\

Tijdens de evaluatie wordt een vereenvoudigde visualisatie van de lineageketen getoond, waarin zichtbaar is uit welke databronnen en tussenliggende tabellen een rapport zijn cijfers haalt en welke rapporten geraakt kunnen worden door wijzigingen in onderliggende datasets. Respondenten wordt gevraagd in welke mate dit soort overzicht hen helpt om cijfers te begrijpen, wijzigingen te verklaren, impact in te schatten en vertrouwen te hebben in rapporten. Voor technische profielen wordt bovendien nagegaan in welke mate de lineage-overzichten hen ondersteunen bij het plannen, analyseren en onderbouwen van wijzigingen.\\

De ingevulde vragenlijsten worden vervolgens geanalyseerd via een combinatie van beschrijvende analyse van de gesloten Likert-vragen en een thematische interpretatie van de open antwoorden. Daarbij wordt gezocht naar terugkerende patronen, zoals ervaren voordelen, knelpunten, ontbrekende informatie en voorwaarden voor breder gebruik van lineage-overzichten.  Op basis hiervan worden conclusies geformuleerd over de evaluatievragen en aanbevelingen geformuleerd in Hoofdstuk \ref{ch:conclusie}\\

\textbf{Deliverable}: een analyse van de gebruikersfeedback op de proof-of-concept, inclusief beantwoording van de evaluatievragen en van de deelvraag over de ervaren meerwaarde van OpenMetadata-lineage, als basis voor de conclusies en aanbevelingen in het slot hoofdstuk.
